
\chapter{TECHNOLOGY STACK}

\section{OVERVIEW}

Stockify uses a modern, pragmatic full-stack chosen for developer productivity, reliability, and scalability.

\section{FRONTEND}

\begin{itemize}
\item React (component-driven UI)
\item TypeScript (static types)
\item Vite (fast dev server and builds)
\item Tailwind CSS (utility-first styling)
\end{itemize}

\section{BACKEND}

\begin{itemize}
\item Node.js (runtime)
\item Express (API framework)
\item Mongoose (MongoDB ODM)
\end{itemize}

\section{DATABASE}

\begin{itemize}
\item MongoDB Atlas — cloud-hosted, scalable document store
\end{itemize}

\section{AUTH \& SECURITY}

\begin{itemize}
\item JWT for stateless auth
\item Google OAuth 2.0 for social login
\item Helmet, CORS and bcrypt used on the backend
\end{itemize}

\section{DEVELOPMENT TOOLS}

\begin{itemize}
\item Editor: VS Code
\item VCS: Git (GitHub)
\item Package manager: npm
\item API testing: Postman
\item Code quality: ESLint, Prettier, TypeScript
\end{itemize}

\section{LIBRARIES}

\begin{itemize}
\item Frontend: React Router, Recharts, React Hook Form, Zod
\item Backend: Joi, bcrypt, CORS, Helmet
\item File/PDF: Tesseract.js, PDFKit, Multer
\end{itemize}

\section{RATIONALE AND CONSIDERATIONS}

Choices favor a single-language stack (JS/TS) for consistency, fast developer feedback (Vite), and flexible data modeling (MongoDB). Performance practices include code-splitting, lazy loading, efficient DB indexing, and caching. Scalability is supported via stateless APIs, Atlas scaling, CDNs, and load balancing.

\section{CONCLUSION}

This stack balances modern developer ergonomics with production-proven components to deliver a maintainable, scalable application.